%~Mouliné par MaN_auto v.0.23.0 2020-06-18 09:32:49
\documentclass[AHL,Unicode,longabstracts]{cedram}

\makeatletter
\def\@altabstractlanguage{}
\makeatother
\usepackage{bbm}
\usepackage{mathrsfs}

%%PACKAGES DESACTIVES
%%\usepackage{etoolbox}
%%\usepackage[retainorgcmds]{IEEEtrantools}
%%%%%%%%%%%%%%%%%%%%%%%%%%%%%


%\DeclareMathOperator{\sgn}{sgn}
%\DeclareMathOperator{\Id}{Id}
%\DeclareMathOperator{\dom}{dom}
%\DeclareMathOperator{\supp}{supp}
%\DeclareMathOperator{\argmin}{argmin}
%\DeclarePairedDelimiter\abs{\lvert}{\rvert}
%\DeclarePairedDelimiter\norm{\lVert}{\rVert}





%%\AtBeginDocument{%\newtheorem{myname}[cdrthm]{My beautiful theorem} }
%\newcommand*\CDRcopyright{License CC BY SA 4.0 https://creativecommons.org/licenses/by-sa/4.0/}
%%\makeatletter\newcommand*\copyright@string{Published under license \href{https://creativecommons.org/licenses/by-sa/4.0/}{CC BY SA 4.0}.}
%\newcommand*\copyright@notice{\copyright@string\\[4pt]\includegraphics[height=18pt]{by-sa.png}}\makeatother


%
%%\providecommand{\email}[1]{\href{mailto:#1}{\nolinkurl{#1}}}
%\providecommand{\arxivref}[2]{\href{http://arxiv.org/abs/#1}{arXiv:#1}\ifblank{#2}{}{ [#2]} }

\newcommand\RR{\mathbb{R}}
\newcommand\Zz{\mathbf{Z}}
\newcommand\Pp{\mathcal{P}}
\newcommand\PP{\mathbb{P}}
\newcommand\EE{\mathbb{E}}
\newcommand\Indic[1]{\Ind_{\{#1\}}}
\newcommand\Ind{\mathbbm{1}}
\newcommand\from{\colon}
\newcommand\NN{\mathbb{N}}
\newcommand\FF{\mathcal{F}}
\newcommand\Uu{\mathcal{U}}
\newcommand\Gg{\mathcal{G}}
\newcommand\Ee{\mathcal{E}}
\newcommand\pp{\mathbf{p}}
\newcommand\e{\mathrm{e}}
\newcommand\Zzu{\Zz^{(u)}}

\newcommand{\crochet}[1]{{\langle #1 \rangle}}

\makeatletter
\newcommand\mathof[1]{{\operator@font#1}} \makeatother
\newcommand\dd{\mathof{d}}

\DeclarePairedDelimiterX\ip[2]{\langle}{\rangle}{#1,#2}

%\newcommand\iu{\mathof{i}}


%environnements à supprimer !! 
%%\newenvironment{eqnarr}{\begin{IEEEeqnarray}{rCl{\ignorespacesafterend}
%%\end{IEEEeqnarray}\ignorespacesafterend}

%%\newenvironment{eqnarr*}{\begin{IEEEeqnarray*}{rCl}}{\end{IEEEeqnarray*}\ignorespacesafterend}
%%\IEEEeqnarraydefcolsep{9}{4em}


%\newcommand{\eqnarrLHS}[1]{\IEEEeqnarraymulticol{3}{l}{#1} \\ \quad}
%\newcommand{\eqnarrLHSnn}[1]{\IEEEeqnarraymulticol{3}{l}{#1} \nonumber \\ \quad}
%\newcommand\dee{\partial}
%\newcommand\Dom{\mathcal{D}}

%\newcommand\nfrac[2]{#1/#2}
%\newcommand\sP{\mathrm{P}}
%\newcommand\sE{\mathrm{E}}
%\newcommand\sQ{\mathrm{Q}}
%\newcommand\downto{\downarrow}
%\newcommand\upto{\uparrow}

%\newcommand\Aa{\mathcal{A}}

%\newcommand\Hh{\mathcal{H}}






%%%%%%%%%%%%%%%%%%%%%%%%%%%%%%%%%%%%%%%%%%%%%%%%%%%%%%%%

\graphicspath{{./figures/}}

\newcommand*{\mk}{\mkern -1mu}
\newcommand*{\Mk}{\mkern -2mu}
\newcommand*{\mK}{\mkern 1mu}
\newcommand*{\MK}{\mkern 2mu}

\hypersetup{urlcolor=purple, linkcolor=blue, citecolor=red}

\newcommand*{\romanenumi}{\renewcommand*{\theenumi}{\roman{enumi}}}
\newcommand*{\Romanenumi}{\renewcommand*{\theenumi}{\Roman{enumi}}}
\newcommand*{\alphenumi}{\renewcommand*{\theenumi}{\alph{enumi}}}
\newcommand*{\Alphenumi}{\renewcommand*{\theenumi}{\Alph{enumi}}}
\let\oldtilde\tilde
\renewcommand*{\tilde}[1]{\mathchoice{\widetilde{#1}}{\widetilde{#1}}{\oldtilde{#1}}{\oldtilde{#1}}}
%%%%%%%%%%%%%%%%%%%%%%%%%%%%%%%%%%%%%%%%%%%%%%%%%%%%%%%%%

